% https://tug.org/texinfohtml/latex2e.html#Font-sizes     see for font sizes!
% \PassOptionsToPackage{showframe}{geometry} % Uncomment to enable showframe
\usepackage[hmargin=1.25in, vmargin=1in, left=3.45cm, right=3.45cm, headheight=21.5764pt, headsep=18pt]{geometry}

\usepackage{fontspec}
\usepackage{polyglossia}
\setmainlanguage{greek}
\setotherlanguage{english}

% windows
\setmainfont{Calibri}[Script=Greek]
\newfontfamily\greekfont{Calibri}[Script=Greek]
\newfontfamily\englishfont{Calibri}
\newfontfamily\greekfonttt{Calibri}[Script=Greek]

% linux
% \setmainfont{Carlito}[Script=Greek, Language=Greek]
% \newfontfamily\greekfont{Carlito}[Script=Greek, Language=Greek]
% \newfontfamily\englishfont{Carlito}
% \newfontfamily\greekfonttt{Carlito}[Script=Greek]

\usepackage{setspace}
\renewcommand{\baselinestretch}{1.5}

\usepackage[table]{xcolor}
\usepackage{float}
\usepackage{placeins}
\usepackage{multirow}
\usepackage{graphicx}
\graphicspath{{images/}}
\usepackage{tikz}
\usepackage{pdflscape}
\usepackage{eso-pic}
\usepackage{caption}
\usepackage{subcaption}
% Figure captions: italic only
\captionsetup[figure]{labelfont={it}, textfont={it}}
% Table captions: bold and italic
\captionsetup[table]{labelfont={bf,it}, textfont={bf,it}, justification=raggedright, singlelinecheck=false}
\usepackage{titlesec}
\definecolor{wordblue}{HTML}{376092}
\definecolor{boxgray}{HTML}{7F7F7F}
\usepackage{enumitem}
\setlist{nosep}

\usepackage{textcomp}
\usepackage[utf8]{inputenc}
\usepackage{booktabs}
\usepackage{makecell}

\usepackage{fancyhdr}
\usepackage{tcolorbox}

\usepackage{longtable}
\usepackage{array}
\usepackage{tablefootnote}
\usepackage[bottom]{footmisc}
\usepackage{amsmath}
\usepackage{amsfonts}
\usepackage{bm}
% The current thesis content does not use \SI; keeping the template
% independent of the optional siunitx package allows compilation on this WSL setup.

\usepackage[backend=biber,style=ieee]{biblatex}
% A document outside the repository root can redefine this before loading
% the preamble (e.g. \newcommand{\bibresourcepath}{../../references.bib}).
\providecommand{\bibresourcepath}{references.bib}
\addbibresource{\bibresourcepath}
\AtBeginBibliography{\sloppy}

\setcounter{secnumdepth}{3}

% for picture annotation
\newcommand{\drawVerticalMarker}[4]{%
  \begin{scope}[x={(image.south east)}, y={(image.north west)}]
    % Define x coordinate for vertical line
    \def\x{#1}  % change this to move the line horizontally
    % Draw vertical dotted line from top (y=#2) to bottom (y=#3)
    \draw[dotted, thick] (\x,#2) -- (\x,#3);
    % Add a label at the bottom
    \node[below] at (\x,#3) {(#4)};
  \end{scope}
}

% Define a reusable command for the page number box
\newcommand{\PageNumberBox}{
    \makebox[0pt][l]{%
        \hspace*{-1.4125in}% move to absolute left edge of physical page
        \begin{tcolorbox}[
            colback=boxgray,
            boxrule=0pt,
            width=3cm,
            height=0.5cm,
            arc=0pt,
            halign=left,
            valign=center
        ]
            \textbf{\hspace{0.5cm}\textcolor{white}{\thepage}}
        \end{tcolorbox}
    }
}

\newcommand{\PageNumberBoxRight}{
    \hspace*{6.875in}% move to absolute right edge of physical page
    \makebox[0pt][r]{%
        \begin{tcolorbox}[
            colback=boxgray,
            boxrule=0pt,
            width=3cm,
            height=0.5cm,
            arc=0pt,
            halign=left,
            valign=center
        ]
            \textbf{\hspace{0.5cm}\textcolor{white}{\thepage}}
        \end{tcolorbox}
    }
}

\renewcommand{\chaptermark}[1]{\markboth{#1}{}}

\fancypagestyle{mystyle}{
    \fancyhf{}

    \fancyhead[L]{%
        \raisebox{3pt}{%
            \ifodd\value{page}
            \else
                \PageNumberBox
            \fi
        }%
    }

    \fancyhead[C]{%
        \raisebox{3pt}{%
            \ifodd\value{page}
            \parbox[t]{\textwidth}{\vspace{-1.4em}\centering\small \thesistitle}
            \else
            \parbox[t]{\textwidth}{\vspace{-1.2em}\centering\small \nouppercase{\leftmark}}
            \fi
        }%
    }

    \fancyhead[R]{%
        \raisebox{3pt}{%
            \ifodd\value{page}
                \PageNumberBoxRight
            \fi
        }%
    }

    \renewcommand{\headrulewidth}{0pt}
}

\usepackage{todonotes}

\fancypagestyle{noboxheader}{
    \fancyhf{}  % Clear all header and footer fields

    \fancyhead[L]{%
        \raisebox{3pt}{%
            \ifodd\value{page}
            \else
                \PageNumberBox
            \fi
        }%
    }

    \renewcommand{\headrulewidth}{0pt}
    \renewcommand{\footrulewidth}{0pt}
}
% Define the first-page-of-chapter style without the center header
\fancypagestyle{plain}{
    \fancyhf{}
    \renewcommand{\headrulewidth}{0pt}
}
\renewcommand{\titlerule}{\rule{\textwidth}{2pt}}

% Apply the custom page style globally
\pagestyle{mystyle}
\makeatletter
\let\ps@plain\ps@plain % Ensure plain style applies to first page of chapters
\makeatother

% Chapter formatting
\titleformat{\chapter}[block]
  {\normalfont\fontsize{20pt}{24pt}\bfseries\itshape\color{wordblue}\raggedleft}  % Right-aligned with raggedleft
  {\chaptername~\thechapter:}  % Chapter label
  {4pt}                        % Space between label and title
  {}                           % Before code
  [\vspace{40pt}\color{wordblue}\titlerule\vspace{1.25cm}]

\newcommand{\MyNumberlessChapter}[1]{%
    % Temporarily redefine numberless chapter style
    \titleformat{name=\chapter,numberless}[block]
      {\normalfont\fontsize{20pt}{24pt}\bfseries\itshape\color{wordblue}\raggedleft}
      {}
      {0pt}
      {}
      [\vspace{40pt}\color{wordblue}\titlerule\vspace{1.25cm}]
    %
    % Manually increment chapter counter and set section numbering prefix
    \refstepcounter{chapter}% increment chapter counter (but don’t print number)
    \chapter*{#1}% print title
    \addcontentsline{toc}{chapter}{#1}% add to TOC
    \markboth{#1}{#1}% update headers
    %
    % Reset section numbering within this “chapter”
    \renewcommand{\thesection}{\Alph{chapter}.\arabic{section}}%
    \setcounter{section}{0}%
}

% Define Unnumbered Chapter Format (\chapter*)
\titleformat{name=\chapter,numberless}[block]
  {\normalfont\Large\bfseries\color{wordblue}}
  {}
  {0pt}
  {\vspace{-1cm}\textsc}

\newcommand{\TopTitleChapter}[2]{%
    % Temporarily redefine numberless chapter style
    \titleformat{name=\chapter,numberless}[block]
      {\normalfont\Large\bfseries\color{wordblue}}
      {}
      {0pt}
      {}%
    % Typeset the heading manually
    {\vspace{5cm}% adjust vertical space if needed
     {\centering\Large\bfseries\color{wordblue}\textsc{#1}\par}% top title
     \vspace{0.5em}%
     {\raggedright\Large\bfseries\color{wordblue}\textsc{#2}\par}% label
     \vspace{1.5em}% <-- space between label and following text
    }%
}

% Section formatting to match the Word document (14pt, bold, 18pt spacing before, and 1.02cm indent)
\titleformat{\section}
  {\fontsize{14pt}{0pt}\selectfont\bfseries} % Set section font to 14pt and bold
  {\thesection}{1em}{}
  \titlespacing*{\section}{0pt}{0.9cm}{0.3cm} % Add spacing before and indentation

% Subsection formatting (optional, if you want it styled)
\titleformat{\subsection}
  {\fontsize{14pt}{0pt}\selectfont\bfseries}
  {\thesubsection}{1em}{}
  \titlespacing*{\subsection}{0pt}{0.8cm}{0.3cm}

% Subsubsection formatting – THIS FIXES YOUR ISSUE
\titleformat{\subsubsection}
  {\fontsize{14pt}{0pt}\selectfont\bfseries}
  {\thesubsubsection}{1em}{}
  \titlespacing*{\subsubsection}{0pt}{0.7cm}{0.3cm}

\newcommand{\unnumberedsection}[1]{%
  \vspace{1em}%
  \noindent\textbf{\fontsize{14pt}{16pt}\selectfont #1}%
  \par\vspace{0.5em}%
}

\numberwithin{equation}{chapter} % Number equations as chapter.equation

\usepackage[
    unicode,
    urlcolor=blue,
    pdfborder={0 0 0},
    pdfauthor={\student},
    pdftitle={\thesistitle}
]{hyperref}

\makeatletter
\renewcommand{\l@table}{\@dottedtocline{1}{1.5em}{3em}}
\renewcommand{\l@figure}{\@dottedtocline{1}{1.5em}{3em}}
\makeatother
